
\def\PUDcodigo{EME110}

\def\PUDcodacad{04506.23}

\def\PUDdisciplina{Materiais II}

\def\PUDsemestre{S5}

\def\PUDhoras{80}

%NOVOS ---------------------------------------------------
\def\PUDhorasTeoria{60}
\def\PUDhorasPratica{20}

\def\PUDinterECA{
\hyperref[PUD:EME108]{Materiais I (04506.20)}

\hyperref[PUD:EME113]{Processos de Fabricação I (04506.24)}

\hyperref[PUD:EME112]{Soldagem (04506.29)}

\hyperref[PUD:EME125]{Análise de Falhas (04506.44)} 
}

\def\PUDatuaisECA{---}

\def\PUDrecursos{Projetor multimídia; Quadro branco e pincel; Aulas em laboratório de materiais.}

%----------------------------------------------------------

\def\PUDcreditos{4}

\def\PUDprerequisito{ \hyperref[PUD:EME108]{EME108} }

\def\PUDpreacad{ \hyperref[PUD:EME108]{Materiais I (04506.20)} }

\def\PUDobjetivos{Conhecer os mecanismos de aumento da resistência em metais;  Conhecer os fenômenos metalúrgicos através dos diagramas de equilíbrio de fases, das curvas TTT e das curvas TRC;  Conhecer os tratamentos térmicos e termoquímicos e suas implicações nos metais;  Conhecer a influência dos elementos de liga nos metais;  Conhecer os processos de fabricação dos aços;  Conhecer os tipos, propriedades e aplicações dos ferros fundidos.}

\def\PUDementa{Discordâncias e Mecanismos de Aumento da Resistência;  Diagramas de Fases;  Transformações de Fases;  Aplicações e Processamento de Ligas Metálicas;  Tratamentos Termoquímicos;  Influência dos Elementos de Liga nos Aços;  Processos de Fabricação dos Aços;  Ferros Fundidos.}

\def\PUDconteudo{ & Discordâncias e deformação plástica.  Sistemas de escorregamento.  Deformação plásticas dos materiais cristalinos.  Mecanismos de aumento da resistência em metais.  Recuperação, recristalização e crescimento de grão. 
\\ 
 & Equilíbrio de fases.  Diagramas de fases binários.  Desenvolvimento de microestruturas.  Sistema Ferro-Carbono.  Desenvolvimento de microestruturas em ligas ferro-carbono.  Influência de outros elementos de liga. 
\\ 
 & Transformações de fases.  Cinètica das transformações de fases.  Alterações microestruturais das ligas ferro-carbono.  Alterações das propriedades em ligas ferro-carbono.  Curvas TTT.  Curvas TRC. 
\\ 
 & Tipos de ligas metálicas.  Fabricação de metais.  Processamento térmico dos metais.  Tratamentos térmicos dos aços.  Recozimento.  Normalização.  Têmpera.  Temperabilidade.  Revenimento.  Martêmpera.  Austêmpera.  Têmpera superficial. 
\\ 
 & Difusão e solubilidade dos elementos químicos.  Cementação.  Nitretação.  Cianetação.  Carbonitretação.  Nitrocarbonetação.  Boretação.  Tratamentos termoreativos.  Microestruturas obtidas. 
\\ 
 & Influência dos elementos de liga nos aços.  Efeito dos elementos de liga na formação da ferrita.  Efeito dos elementos de liga na formação da perlita.  Efeitos dos elementos de liga nos carbonetos, nas inclusões não-metálicas e nos compostos intermetálicos.  Efeito dos elementos de liga na têmpera e no revenimento.  Efeito dos principais elementos de liga nos aços.  Impurezas dos aços. 
\\ 
 & Processos de fabricação dos aços.  Produção de ferro-gusa.  Processos de redução direta.  Aciaria.  Lingotamento e lingotes.  Processos especiais de refino e otenção de aços e ligas especiais. 
\\ 
 & Ferros fundidos.  Tipos de ferros fundidos. }

\def\PUDmetodo{Aulas expositórias que podem ser teóricas e/ou práticas, onde as práticas no laboratório serão marcadas no decorrer da disciplina.}

\def\PUDavalia{A avaliação será feita com aplicação de provas teóricas e/ou práticas, além da possibilidade de inclusão de trabalhos e seminários no decorrer da disciplina.}

\def\PUDbibbasica{ & \href{www.biblioteca.ifce.edu.br/index.asp?codigo_sophia=24421}{Callister Jr.}, W. D.; Rethwisch, D. G.  Ciência e Engenharia de Materiais - Uma Introdução.  8ª ed.  Rio de Janeiro, RJ: LTC, 2012.  817p.  ISBN: 9788521621249. 
\\ 
 & \href{www.biblioteca.ifce.edu.br/index.asp?codigo_sophia=24424}{Silva}, André Luiz V. da Costa e.  Aços e Ligas Especiais.  3ª ed.  São Paulo, SP: Blucher, 2010.  646p.  ISBN: 9788521205180. 
\\ 
 &  \href{www.biblioteca.ifce.edu.br/index.asp?codigo_sophia=39995}{Chiaverini}, V.  Aços e ferros fundidos: características gerais, tratamentos térmicos e principais tipos.  7ª ed.  São Paulo, SP: ABM, 2012.  599p.  ISBN: 8586778486. }

\def\PUDbibcomplementar{ & \href{www.biblioteca.ifce.edu.br/index.asp?codigo_sophia=42347}{Shackelford}, J. F.  Ciência dos Materiais.  6ª ed.  São Paulo, SP: Pearson, 2012.  556p.  ISBN: 9788576051602. 
\\ 
 & \href{www.biblioteca.ifce.edu.br/index.asp?codigo_sophia=23980}{Chiaverini}, V.  Tecnologia Mecânica: Estrutura e Propriedades das Ligas Metálicas.  2ª ed.  São Paulo, SP: Pearson Education do Brasil, 1986.  266p.  ISBN: 0074500899. 
\\ 
 & \href{www.biblioteca.ifce.edu.br/index.asp?codigo_sophia=23701}{Chiaverini}, V.  Tecnologia Mecânica: Materiais de Construção Mecânica.  v. 3.  2ª ed.  São Paulo, SP: Pearson Education do Brasil, 1986.  388p.  ISBN: 9780074500910. 
\\
 & \href{biblioteca.ifce.edu.br/index.asp?codigo_sophia=24118}{Colpaert}, Hubertus. Metalografia dos produtos siderúrgicos comuns. 4. ed. São Paulo: Edgard Blücher, 2008. 652 p. + Inclui CD-ROM. Inclui bibliografia e índice. ISBN 9788521204497.
\\
 & \href{www.biblioteca.ifce.edu.br/index.asp?codigo_sophia=58235}{Pavanati}, Henrique Cezar.  Ciência e Tecnologia dos Materiais.  E-book.  Pearson.  196p.  ISBN: 9788543009797.
%\\ 
%& \href{biblioteca.ifce.edu.br/index.asp?codigo_sophia=47009}{Nunes}, Laerce de Paula. Materiais: aplicações de engenharia, seleção e integridade. Rio de Janeiro: Interciência, 2012. 375 p. ISBN 9788571932883.
}

\def\PUDautor{Francisco Nélio Costa Freitas}

\def\PUDdata{2013-05-20}

\def\PUDrevisao{2}

\def\PUDrevisor{Venceslau Xavier de Lima Filho}

\def\PUDdatarevisao{2019-05-15}

\PUDcorpo
